% ==========================================================
%  OLD EXAM QUESTIONS — Unsupervised Machine Learning
%  (WORKED SOLUTIONS)
%  Companion to exams.tex: every question is followed
%  immediately by a green answer box. For multiple-choice,
%  correct options are marked with a green check, incorrect
%  ones with a red cross, each with a short justification
%  (incl. why tricky near-miss options are wrong). For
%  calculations the working is shown; wild distractors that
%  only exist to catch guessers get no explanation.
%
%  NOTE: exams.tex stays the blank practice set. This file is
%  the solutions variant — do not merge the two.
% ==========================================================

\documentclass[11pt,oneside]{article}

\usepackage[
  a4paper,
  top=2.4cm,
  bottom=2.1cm,
  left=2.3cm,
  right=2.3cm,
  headheight=15pt
]{geometry}

\usepackage[T1]{fontenc}
\usepackage[utf8]{inputenc}
\usepackage{lmodern}
\usepackage{microtype}
\usepackage{inconsolata}
\usepackage{amsmath,amssymb,mathtools,bm}
\usepackage{enumitem}
\usepackage{booktabs,array,multicol,longtable}
\usepackage{xcolor}
\usepackage[most]{tcolorbox}
\usepackage{titlesec}
\usepackage{fancyhdr}
\usepackage[hidelinks]{hyperref}

\setlength{\parindent}{0pt}
\setlength{\parskip}{0.55em}
\setlist[itemize]{topsep=3pt,itemsep=2pt,leftmargin=2.2em}
\setlist[enumerate]{topsep=3pt,itemsep=2pt,leftmargin=2.2em}

% ── Colors ───────────────────────────────────────────────
% Shared accent with summary.tex — do not change.
\definecolor{accent}{HTML}{1B3A5C}
\definecolor{q@frame}{HTML}{1A5276}
\definecolor{q@bg}{HTML}{F6FAFD}
\definecolor{key@frame}{HTML}{1E8449}
\definecolor{key@bg}{HTML}{F4FBF6}
\definecolor{note@frame}{HTML}{9A6A00}
\definecolor{note@bg}{HTML}{FEF9E7}
\definecolor{ok}{HTML}{1E8449}    % green check — correct option
\definecolor{bad}{HTML}{B03A2E}   % red cross  — incorrect option

% ── Box global style ─────────────────────────────────────
\tcbset{
  enhanced,
  breakable,
  arc=2.5mm,
  boxrule=0.5pt,
  left=4.5mm,
  right=4mm,
  top=2.5mm,
  bottom=2.5mm,
  fonttitle=\bfseries\small,
  sharp corners=west,
  parbox=false
}

\newtcolorbox{qbox}[1][]{
  colback=q@bg,
  colframe=q@frame,
  borderline west={3.5pt}{0pt}{q@frame},
  title={#1}
}

\newtcolorbox{keybox}[1][]{
  colback=key@bg,
  colframe=key@frame,
  borderline west={3.5pt}{0pt}{key@frame},
  title={#1}
}

\newtcolorbox{notebox}[1][]{
  colback=note@bg,
  colframe=note@frame,
  borderline west={3.5pt}{0pt}{note@frame},
  title={#1}
}

% ── Section headings ─────────────────────────────────────
\titleformat{\section}
  {\normalfont\Large\bfseries\color{accent}}
  {}{0em}{}
  [{\vspace{1pt}\color{accent!45}\titlerule[0.45pt]}]
\titlespacing*{\section}{0pt}{3.2ex plus 1ex minus .2ex}{1.8ex plus .2ex}

\titleformat{\subsection}
  {\normalfont\normalsize\bfseries\color{accent!80!black}}
  {}{0em}{}
\titlespacing*{\subsection}{0pt}{2.2ex plus 1ex minus .2ex}{1.0ex plus .2ex}

% ── Header & footer ──────────────────────────────────────
\pagestyle{fancy}
\fancyhf{}
\fancyhead[L]{\small\color{gray!70}Old Exam Questions (Solutions) --- \coursetitle}
\fancyhead[R]{\small\color{gray!70}\thepage}
\renewcommand{\headrulewidth}{0.4pt}

% ── Document metadata ─────────────────────────────────────
\newcommand{\coursetitle}{Unsupervised Machine Learning}

% ── Question counter & macros ────────────────────────────
\newcounter{question}
\newcommand{\questiontitle}[1]{%
  \stepcounter{question}\begin{qbox}[Question \thequestion: #1]}
\newcommand{\closequestion}{\end{qbox}}
\newcommand{\repeatnote}{\textit{\color{gray!70}Asked multiple times.}}
\newcommand{\choice}{\item[$\square$] }
\newcommand{\answerline}[1]{\makebox[#1]{\hrulefill}}

% ── Inline answer macros (solutions variant) ─────────────
% Green answer box placed directly under each question.
% Multiple-choice usage:
%   \begin{answer}{A, B}
%     \yes{A} reason why correct ...
%     \no{C}  reason why wrong (esp. for tricky near-misses) ...
%   \end{answer}
\newenvironment{answer}[1]{%
  \begin{keybox}[Answer]%
  \begin{itemize}[leftmargin=1.7em,itemsep=1.5pt,topsep=2pt]}%
  {\end{itemize}\end{keybox}}
\newcommand{\yes}[1]{\item[\textcolor{ok}{\ensuremath{\checkmark}}]\textbf{#1.}\ }
\newcommand{\no}[1]{\item[\textcolor{bad}{\ensuremath{\bm\times}}]\textbf{#1.}\ }

% Prose answer box — for calculations and labelling questions.
% Usage:  \begin{answerprose}{result}  working / explanation  \end{answerprose}
\newenvironment{answerprose}[1]{%
  \begin{keybox}[Answer]}%
  {\end{keybox}}

% ── Math macros ──────────────────────────────────────────
\newcommand{\R}{\mathbb{R}}
\newcommand{\E}{\mathbb{E}}
\newcommand{\Prob}{\mathbb{P}}
\newcommand{\KL}[2]{D_{\mathrm{KL}}\!\left(#1\,\|\,#2\right)}
\newcommand{\dd}{\,\mathrm{d}}
\newcommand{\Var}{\mathrm{Var}}


% ==========================================================
%  BEGIN DOCUMENT
% ==========================================================
\begin{document}

% ── Title page ───────────────────────────────────────────
\thispagestyle{empty}
\vspace*{1.4cm}
\begin{center}
  {\color{accent}\rule{\linewidth}{1.8pt}}\par
  \vspace{0.65cm}
  {\fontsize{26}{32}\selectfont\bfseries\color{accent}Old Exam Questions\\[2pt] \coursetitle\par}
  \vspace{0.25cm}
  {\Large\color{key@frame}\bfseries Worked Solutions\par}
  \vspace{0.55cm}
  {\color{accent!35}\rule{\linewidth}{0.5pt}}\par
  \vspace{0.55cm}
  {\large Stefan Eder\par}
  \vspace{0.15cm}
  {\small\color{gray!70}\today\par}
\end{center}

\vspace{0.9cm}
\begin{notebox}[How this file is organized]
This is the \textbf{solutions companion} to \texttt{exams.tex}. Every question is followed
\emph{immediately} by a green \textbf{answer box}. For \textbf{multiple-choice} questions each
option is marked \textcolor{ok}{\ensuremath{\checkmark}} (correct) or
\textcolor{bad}{\ensuremath{\bm\times}} (incorrect), with a short justification --- \emph{why} the
correct options are correct and, for the tricky / near-miss distractors, \emph{why} they are wrong.
For \textbf{calculations} the working is shown. Options are referred to by position (A = first
option, B = second, \dots). If you want a clean practice set with blank answers, use
\texttt{exams.tex}.
\end{notebox}

\tableofcontents
\newpage

\section{Probability, Information Theory, and Estimation}

\questiontitle{Cumulants for independent uniform variables}
Assume you have two independent random variables $X$ and $Y$ which are uniformly distributed on $[-1,1]$.
As in the lecture, let $\kappa_n(X)$ denote the $n$-th order cumulant of a random variable $X$.
Compute the following quantities and answer the last classification question.
Use the hint
\[
\kappa_4(X)=\E(X^4)-3\bigl(\E(X^2)\bigr)^2
\]
for centered $X$.

\begin{itemize}
  \item $\kappa_4(X)$
  \item $\kappa_4(2X)$
  \item $\kappa_4(X+Y)$
  \item $X$ is sub-Gaussian / super-Gaussian
\end{itemize}
\closequestion
\begin{answerprose}{-0.13 / -2.13 / -0.27 / sub-Gaussian}
For $X\sim U[-1,1]$: $\E(X^2)=\tfrac13$ and $\E(X^4)=\tfrac15$ (in general $\E(X^{2k})=\tfrac{1}{2k+1}$).
\[
\kappa_4(X)=\tfrac15-3\bigl(\tfrac13\bigr)^2=\tfrac15-\tfrac13=-\tfrac{2}{15}\approx\mathbf{-0.13}.
\]
Cumulants are homogeneous of degree $n$: $\kappa_4(cX)=c^4\kappa_4(X)$, so
$\kappa_4(2X)=16\cdot(-\tfrac{2}{15})=-\tfrac{32}{15}\approx\mathbf{-2.13}$.
Cumulants are additive for independent variables:
$\kappa_4(X+Y)=\kappa_4(X)+\kappa_4(Y)=-\tfrac{4}{15}\approx\mathbf{-0.27}$.
Since $\kappa_4(X)<0$ (negative excess kurtosis), $X$ is \textbf{sub-Gaussian}
--- the uniform distribution has lighter tails than a Gaussian.
\end{answerprose}

\questiontitle{Cumulants for symmetric Bernoulli variables}
Assume you have two independent, discrete random variables $X$ and $Y$, both following independent symmetric Bernoulli distributions taking only the values $+1$ and $-1$, each with probability $1/2$.
As in the lecture, let $\kappa_n(X)$ denote the $n$-th order cumulant of a random variable $X$.
Compute the quantities below.
Use the hint
\[
\kappa_3(X)=\E(X^3), \qquad \kappa_4(X)=\E(X^4)-3\bigl(\E(X^2)\bigr)^2
\]
for centered $X$.

\begin{itemize}
  \item $\kappa_3(X)$
  \item $\kappa_4(X)$
  \item $\kappa_4(3X)/\kappa_4(X)$
  \item $\kappa_4(X+Y)/\kappa_4(X)$
\end{itemize}
\closequestion
\begin{answerprose}{0 / -2 / 81 / 2}
$X^2=X^4=1$ always, and $X^3=X$, so
\[
\kappa_3(X)=\E(X^3)=\E(X)=\mathbf{0}
\qquad\text{(the distribution is symmetric),}
\]
\[
\kappa_4(X)=\E(X^4)-3\bigl(\E(X^2)\bigr)^2=1-3=\mathbf{-2}.
\]
Homogeneity: $\kappa_4(3X)/\kappa_4(X)=3^4=\mathbf{81}$.
Additivity under independence: $\kappa_4(X+Y)=2\,\kappa_4(X)$, so the ratio is $\mathbf{2}$.
\end{answerprose}

\questiontitle{Eigenvectors, covariance matrix, and a scalar expression}
You are given a dataset with $n$ samples and three features.
The covariance matrix $C$ has eigenvalues $(\lambda_1,\lambda_2,\lambda_3)=(5,3,1)$
and associated normalized eigenvectors $(v_1,v_2,v_3)$ with $\lVert v_i\rVert=1$.
Compute
\[
A=(2v_1^\top+v_2^\top+3v_3^\top)\,C\,(3v_2+v_3).
\]
\closequestion
\begin{answerprose}{12}
Use $Cv_i=\lambda_i v_i$ and orthonormality of the eigenvectors ($v_i^\top v_j=\delta_{ij}$):
\[
C(3v_2+v_3)=3\lambda_2 v_2+\lambda_3 v_3=9v_2+v_3.
\]
Multiplying from the left kills every mixed term:
\[
A=(2v_1^\top+v_2^\top+3v_3^\top)(9v_2+v_3)
=9\,(v_2^\top v_2)+3\,(v_3^\top v_3)=9+3=\mathbf{12}.
\]
\end{answerprose}

\questiontitle{KL divergence for continuous uniform distributions}
\repeatnote

Assume $Y$ and $Z$ are uniformly distributed on $[0,\theta_Y]$ and $[0,\theta_Z]$, respectively:
\[
p_Y(x)=\begin{cases}
\frac{1}{\theta_Y}, & x\in[0,\theta_Y],\\
0, & \text{else,}
\end{cases}
\qquad
p_Z(x)=\begin{cases}
\frac{1}{\theta_Z}, & x\in[0,\theta_Z],\\
0, & \text{else.}
\end{cases}
\]
For $\theta_Y=e^5$ and $\theta_Z=e^{10}$, compute
\[
D_{\mathrm{KL}}(p_Y\|p_Z)=\int_{-\infty}^{\infty} p_Y(x)\log\!\left(\frac{p_Y(x)}{p_Z(x)}\right)\dd x.
\]
Use the natural logarithm and round to two decimal places.
\closequestion
\begin{answerprose}{5.00}
The integrand is nonzero only on $[0,\theta_Y]$ (where $p_Y>0$), and there the ratio is constant:
\[
D_{\mathrm{KL}}(p_Y\|p_Z)
=\int_0^{\theta_Y}\frac{1}{\theta_Y}\ln\!\frac{1/\theta_Y}{1/\theta_Z}\dd x
=\ln\frac{\theta_Z}{\theta_Y}
=\ln\frac{e^{10}}{e^{5}}=10-5=\mathbf{5.00}.
\]
(This works because the support of $p_Y$ is contained in the support of $p_Z$; the reverse
divergence would be infinite.)
\end{answerprose}

\questiontitle{Outer product matrix}
Given two vectors
\[
x=\begin{pmatrix}1\\2\\3\end{pmatrix},
\qquad
y=\begin{pmatrix}4\\5\\6\end{pmatrix},
\]
compute the outer product matrix $A=xy^\top$ and give the third row of $A$.
\closequestion
\begin{answerprose}{(12, 15, 18)}
$(xy^\top)_{ij}=x_i y_j$, so row $i$ of $A$ is $x_i\cdot y^\top$. The third row is
\[
3\cdot(4,5,6)=(\mathbf{12},\,\mathbf{15},\,\mathbf{18}).
\]
\end{answerprose}

\questiontitle{Marginal probabilities, entropies, and mutual information --- variant 1}
\repeatnote

Assume you are given two binary random variables $X$ and $Y$ taking only the values $0$ and $1$ where their joint probability distribution is specified by
\[
\begin{array}{c|cc}
 & Y=0 & Y=1\\
\hline
X=0 & 0.25 & 0.125\\
X=1 & 0.25 & 0.375
\end{array}
\]
Compute $P(X=1)$, $P(Y=1)$, $H(X)$, $H(Y)$, $I(X;Y)$ using the natural logarithm.
Round the final values to two decimal places.
\closequestion
\begin{answerprose}{0.63 / 0.50 / 0.66 / 0.69 / 0.03}
Marginals (row/column sums):
\[
P(X=1)=0.25+0.375=\mathbf{0.63},\qquad
P(Y=1)=0.125+0.375=\mathbf{0.50}.
\]
Entropies:
\[
H(X)=-(0.375\ln 0.375+0.625\ln 0.625)\approx\mathbf{0.66},\qquad
H(Y)=\ln 2\approx\mathbf{0.69}.
\]
Joint entropy:
$H(X,Y)=-(0.25\ln 0.25+0.125\ln 0.125+0.25\ln 0.25+0.375\ln 0.375)\approx 1.32$.
\[
I(X;Y)=H(X)+H(Y)-H(X,Y)\approx 0.662+0.693-1.321=\mathbf{0.03}.
\]
\end{answerprose}

\questiontitle{Marginal probabilities, entropies, and mutual information --- variant 2}
\repeatnote

Assume you are given two binary random variables $X$ and $Y$ taking only the values $0$ and $1$ where their joint probability distribution is specified by
\[
\begin{array}{c|cc}
 & Y=0 & Y=1\\
\hline
X=0 & \tfrac13 & \tfrac13\\
X=1 & 0 & \tfrac13
\end{array}
\]
Compute $P(Y=0)$, $H(X)$, $H(Y)$, $I(X;Y)$ using the natural logarithm.
Round the final values to two decimal places.
\closequestion
\begin{answerprose}{0.33 / 0.64 / 0.64 / 0.17}
Marginals: $P(Y=0)=\tfrac13\approx\mathbf{0.33}$, $P(X=0)=\tfrac23$, $P(Y=1)=\tfrac23$.
\[
H(X)=H(Y)=-\bigl(\tfrac23\ln\tfrac23+\tfrac13\ln\tfrac13\bigr)\approx\mathbf{0.64}.
\]
The joint distribution has three equally likely cells, so $H(X,Y)=\ln 3\approx 1.10$.
\[
I(X;Y)=H(X)+H(Y)-H(X,Y)\approx 0.637+0.637-1.099=\mathbf{0.17}.
\]
\end{answerprose}

\questiontitle{Marginal probabilities, entropies, and mutual information --- variant 3}
\repeatnote

Assume you are given two binary random variables $X$ and $Y$ taking only the values $0$ and $1$ where their joint probability distribution is specified by
\[
\begin{array}{c|cc}
 & Y=0 & Y=1\\
\hline
X=0 & \tfrac{3}{10} & \tfrac{4}{10}\\
X=1 & \tfrac{2}{10} & \tfrac{1}{10}
\end{array}
\]
Compute $P(X=0)$, $P(Y=0)$, $H(X)$, $H(Y)$, $I(X;Y)$ using the natural logarithm.
Round the final values to two decimal places.
\closequestion
\begin{answerprose}{0.70 / 0.50 / 0.61 / 0.69 / 0.02}
Marginals:
\[
P(X=0)=0.3+0.4=\mathbf{0.70},\qquad
P(Y=0)=0.3+0.2=\mathbf{0.50}.
\]
Entropies:
\[
H(X)=-(0.7\ln 0.7+0.3\ln 0.3)\approx\mathbf{0.61},\qquad
H(Y)=\ln 2\approx\mathbf{0.69}.
\]
Joint entropy:
$H(X,Y)=-(0.3\ln 0.3+0.4\ln 0.4+0.2\ln 0.2+0.1\ln 0.1)\approx 1.28$.
\[
I(X;Y)=H(X)+H(Y)-H(X,Y)\approx 0.611+0.693-1.280=\mathbf{0.02}.
\]
\end{answerprose}
\questiontitle{Parameter Estimation and Maximum Likelihood}
\repeatnote

Which statements about Parameter Estimation and Maximum Likelihood are true?
Multiple answers are possible.
\begin{itemize}
  \choice If the data are distributed i.i.d., the likelihood can be written as a product of individual probabilities over the data samples.
  \choice The likelihood and the log-likelihood are maximized for the same parameter value, and at that parameter value $\hat w$, their values coincide: $\mathcal{L}(\hat w)=\ln \mathcal{L}(\hat w)$.
  \choice Parameter estimation and the maximum likelihood approach rest on the assumption that the model class is not known and estimated from the data.
  \choice The maximum likelihood estimator is asymptotically unbiased and efficient, but not asymptotically consistent.
  \choice The mean squared error can be written as the sum of the variance of the estimator and the squared bias of the estimator.
  \choice An estimator $\hat w$ is unbiased if its expected value equals the true parameter: $\E(\hat w)=w$.
  \choice The likelihood $\mathcal{L}(\{x\};w)$ is the probability that the model with parameter $w$ produces the data $\{x\}$.
  \choice The mean squared error is defined as $\mathrm{mse}(\hat w)=(\hat w-w)^\top(\hat w-w)$.
\end{itemize}
\closequestion
\begin{answer}{A, E, F, G}
  \yes{A} Independence lets the joint density factorize: $\mathcal{L}=\prod_i p(x_i;w)$.
  \no{B} Same \emph{argmax} (the logarithm is strictly monotone), but the \emph{values} differ: $\ln\mathcal{L}(\hat w)\neq\mathcal{L}(\hat w)$ in general.
  \no{C} The model class is assumed to be \emph{known and fixed}; only its parameters are estimated from the data.
  \no{D} ML \emph{is} asymptotically consistent (and unbiased and efficient) --- the ``not consistent'' spoils the statement.
  \yes{E} Bias--variance decomposition: $\mathrm{mse}(\hat w)=\Var(\hat w)+\mathrm{bias}^2(\hat w)$.
  \yes{F} That is the definition of unbiasedness.
  \yes{G} That is the definition of the likelihood.
  \no{H} The expectation is missing: $\mathrm{mse}(\hat w)=\E\bigl[(\hat w-w)^\top(\hat w-w)\bigr]$; without it, the expression is a random variable, not an error measure.
\end{answer}

\questiontitle{Dropdown selection task on estimation theory}
\repeatnote

Fill in the correct choices.

For any parameter $w$, an estimator $\hat w$ is \textit{consistent / efficient / unbiased} if $\E(\hat w)=w$.
The \textit{expected value / bias / mean squared error / variance} of $\hat w$ can be decomposed as a sum of its
\textit{expected value / bias / mean squared error / variance} and its squared
\textit{expected value / bias / mean squared error / variance}.

Let $X_1,\dots,X_n$ be random variables drawn independently from a distribution $p(\cdot;w)$.
Consider
\[
\frac{1}{n+1}\sum_{i=1}^n X_i
\]
as an estimator for the mean of the distribution $p(\cdot;w)$.
This estimator is
\textit{neither unbiased nor consistent / unbiased but not consistent / biased but consistent / both unbiased and consistent}.

Now consider $X_1$ as an estimator for the mean of the distribution $p(\cdot;w)$.
This estimator is
\textit{neither unbiased nor consistent / unbiased but not consistent / biased but consistent / both unbiased and consistent}.

Let $\{x_1,\dots,x_n\}$ be a sample drawn independently from the distribution $p(\cdot;w)$.
The function
\[
\prod_{i=1}^n p(x_i;w)
\]
is called the \textit{probability / loss / density / likelihood} function.
Since $p(\cdot;w)$ is usually a \textit{probability / loss / density / likelihood} function, the probability measure of the set $\{x_1,\dots,x_n\}$ is in fact zero.

The Cram\'er--Rao lower bound relates the \textit{expected value / bias / mean squared error / variance} of the estimator $\hat w$ to the inverse of its Fisher information.
An estimator that reaches the Cram\'er--Rao lower bound is said to be \textit{consistent / efficient / unbiased}.
\closequestion
\begin{answerprose}{unbiased / mse / variance / bias / biased but consistent / unbiased but not consistent / likelihood / density / variance / efficient}
In order of the blanks:
\begin{itemize}[leftmargin=1.7em,itemsep=1.5pt,topsep=2pt]
  \item \textbf{unbiased} --- $\E(\hat w)=w$ is the definition of unbiasedness.
  \item \textbf{mean squared error}, decomposed into \textbf{variance} $+$ squared \textbf{bias}:
        $\mathrm{mse}(\hat w)=\Var(\hat w)+\mathrm{bias}^2(\hat w)$.
  \item $\frac{1}{n+1}\sum_i X_i$ is \textbf{biased but consistent}: its expectation is
        $\frac{n}{n+1}\mu\neq\mu$ for finite $n$, but bias and variance both vanish as $n\to\infty$.
  \item $X_1$ is \textbf{unbiased but not consistent}: $\E(X_1)=\mu$, but the estimator never
        concentrates --- it ignores all further data.
  \item $\prod_i p(x_i;w)$ is the \textbf{likelihood} function.
  \item $p(\cdot;w)$ is usually a \textbf{density}, so any finite set of points has probability
        measure zero.
  \item Cram\'er--Rao bounds the \textbf{variance}: $\Var(\hat w)\ge 1/\mathrm{Fisher}$;
        an estimator attaining the bound is \textbf{efficient}.
\end{itemize}
\end{answerprose}

\questiontitle{Unsupervised Machine Learning}
Which statements about Unsupervised Machine Learning are true?
Multiple answers are possible.
\begin{itemize}
  \choice Underfitting means that a model is adapted to special training characteristics such as noisy measurements or outliers.
  \choice Projection methods project data into a higher-dimensional space while keeping their neighborhoods.
  \choice The model selection procedure is called ``training'' or ``learning''.
  \choice Examples of unsupervised methods are: PCA, ICA, Factor Analysis, k-means clustering, and mixture models.
  \choice In parametric models, every parameter (vector) realization represents a model.
  \choice $k$-nearest neighbors is a nonparametric model.
  \choice Generative models do not aim at modeling the data creation process.
  \choice Unsupervised Machine Learning algorithms learn patterns from unlabeled data.
\end{itemize}
\closequestion
\begin{answer}{D, E, F, H}
  \no{A} Adapting to noise and outliers of the training set is \emph{over}fitting; underfitting means the model is too simple to capture the structure at all.
  \no{B} Projection methods map into a \emph{lower}-dimensional space.
  \no{C} Training/learning fits the parameters of a model \emph{within} a model class; model selection is the separate step of choosing among models.
  \yes{D} All of these work without labels --- classic unsupervised toolbox.
  \yes{E} A parametric model class is a family indexed by the parameter vector; each realization picks one model.
  \yes{F} $k$-NN keeps the training data itself instead of a fixed-size parameter vector --- nonparametric.
  \no{G} The opposite: generative models explicitly model the data-generating process.
  \yes{H} That is the definition of unsupervised learning.
\end{answer}

\questiontitle{Expectation Maximization (EM) and Maximum Entropy --- version A}
Which statements about Expectation Maximization (EM) and Maximum Entropy are correct?
Multiple answers are possible.
\begin{itemize}
  \choice The EM algorithm can be applied to hidden Markov models and mixture models.
  \choice In the M-step, the model parameters are updated.
  \choice The E-step increases the upper bound on the log-likelihood.
  \choice In EM, one deals with problems that have hidden variables as well as model parameters.
  \choice The entropy of a probability distribution is always larger equal $0$ and smaller equal $1$.
  \choice Hidden state models in general have analytical solutions for the likelihood.
  \choice In the E-step, the hidden variables parameters are updated.
  \choice The Maximum Entropy approach requires minimal prior assumptions to assign a-priori probabilities.
\end{itemize}
\closequestion
\begin{answer}{A, B, D, G, H}
  \yes{A} Baum--Welch (HMMs) and mixture-model fitting are \emph{the} classic EM applications.
  \yes{B} The M-step maximizes the expected complete-data log-likelihood with respect to the parameters.
  \no{C} EM works with a \emph{lower} bound on the log-likelihood; there is no upper bound involved.
  \yes{D} Exactly the EM setting: latent variables plus model parameters.
  \no{E} Discrete entropy is $\ge 0$ but \emph{not} bounded by $1$: a uniform distribution over $k$ outcomes has entropy $\ln k>1$ for $k\ge 3$.
  \no{F} Latent-variable likelihoods are generally analytically intractable --- that is precisely why EM exists.
  \yes{G} The E-step (re-)estimates the hidden variables via their posterior given the current parameters.
  \yes{H} MaxEnt chooses the least-informative distribution consistent with the known constraints.
\end{answer}

\questiontitle{Expectation Maximization (EM) and Maximum Entropy --- version B}
\repeatnote

Which statements about Expectation Maximization (EM) and Maximum Entropy are correct?
Multiple answers are possible.
\begin{itemize}
  \choice In the M-step, the hidden variables are updated.
  \choice Maximum Entropy requires minimal prior assumptions to assign a-priori probabilities.
  \choice The E-step increases the upper bound on the log-likelihood.
  \choice In the E-step, the model parameters are updated.
  \choice The EM algorithm can be applied to hidden Markov models and mixture models.
  \choice The entropy of a probability distribution is always larger equal $0$ and smaller equal $1$.
  \choice The E-step increases the lower bound on the log-likelihood.
  \choice In EM, one deals with problems that have hidden variables as well as model parameters.
\end{itemize}
\closequestion
\begin{answer}{B, E, G, H}
  \no{A} Swapped: the M-step updates the \emph{model parameters}.
  \yes{B} MaxEnt = minimal prior assumptions.
  \no{C} Lower bound, not upper bound.
  \no{D} Swapped: the E-step updates the estimates of the \emph{hidden variables}.
  \yes{E} HMMs (Baum--Welch) and mixture models are the standard EM examples.
  \no{F} Entropy can exceed $1$ (uniform over $k\ge3$ outcomes: $\ln k>1$).
  \yes{G} The E-step tightens the lower bound so that it touches the log-likelihood at the current parameters.
  \yes{H} That is the EM setting.
\end{answer}
\section{PCA, Kernel PCA, and ICA}

\questiontitle{Principal Component Analysis --- statements, version A}
Which statements about Principal Component Analysis (PCA) are correct?
Multiple answers are possible.
\begin{itemize}
  \choice The first principal component corresponds to the largest eigenvalue in the eigendecomposition of the sample covariance matrix.
  \choice The singular values from the singular value decomposition of the data matrix $X$ are equal to the eigenvalues of the eigendecomposition of $X^\top X$.
  \choice The PCA solution is unique, i.e. there is only one solution.
  \choice Eigenvectors can be computed iteratively with Oja's rule.
  \choice The sample variance of the $k$-th projection is equal to the $k$-th eigenvalue of the sample covariance matrix.
  \choice The matrix of eigenvectors $U$ is an orthogonal matrix.
  \choice The sample covariance matrix is symmetric and positive semidefinite.
  \choice PCA transforms data into a new coordinate system such that the data has the largest mean value along the first coordinate, the second largest mean value along the second coordinate, and so on.
\end{itemize}
\closequestion
\begin{answer}{A, D, E, F, G}
  \yes{A} Principal components are sorted by eigenvalue; the first one carries the largest.
  \no{B} Near-miss: the singular values are the \emph{square roots} of the eigenvalues of $X^\top X$: $\sigma_i=\sqrt{\lambda_i}$.
  \no{C} Not unique --- every eigenvector can be sign-flipped, and with equal eigenvalues the directions are not even determined.
  \yes{D} Oja's rule is the classic iterative (Hebbian) method to extract the leading eigenvector.
  \yes{E} The variance of the data projected onto $u_k$ is exactly $\lambda_k$.
  \yes{F} Eigenvectors of a symmetric matrix can be chosen orthonormal, so $U^\top U=I$.
  \yes{G} Any sample covariance matrix is symmetric and positive semidefinite by construction.
  \no{H} PCA orders directions by \emph{variance}, not by mean value.
\end{answer}

\questiontitle{Principal Component Analysis --- statements, version B}
We keep the notation of the lecture slides: $X$ is the data matrix and $U$ is the matrix of eigenvectors of $X^\top X$.
Which statements about PCA are true?
Multiple answers are possible.
\begin{itemize}
  \choice PCA transforms data into a new coordinate system such that the data has the largest variance along the first coordinate, the second largest variance along the second coordinate, and so on.
  \choice The first principal component corresponds to the largest eigenvalue in the eigendecomposition of the sample covariance matrix.
  \choice If there are $n$ data samples, each of which is $d$-dimensional, the matrix $U$ has dimension $n\times d$.
  \choice The sample covariance matrix is orthogonal and positive semidefinite.
  \choice The singular values from the singular value decomposition of the data matrix $X$ are equal to the square root of the eigenvalues of the eigendecomposition of $X^\top X$.
  \choice The eigenvalues obtained from the eigendecomposition of $X^\top X$ are all positive and sum up to $1$.
  \choice PCA is commonly used to visualize data by downprojecting it to a smaller number of dimensions, keeping the most relevant information.
  \choice The matrix $U$ is a symmetric matrix.
\end{itemize}
\closequestion
\begin{answer}{A, B, E, G}
  \yes{A} That is the variance-maximization view of PCA.
  \yes{B} Correct, as in version A.
  \no{C} $U$ collects the eigenvectors of the $d\times d$ matrix $X^\top X$, so it is $d\times d$ --- the sample count $n$ does not appear.
  \no{D} Symmetric and PSD, yes; \emph{orthogonal}, no --- a covariance matrix is generally not orthogonal.
  \yes{E} $\sigma_i=\sqrt{\lambda_i}$ --- this is the correct version of the near-miss in version A.
  \no{F} Eigenvalues are only \emph{non-negative} (zero is possible), and they sum to the total variance ($\mathrm{tr}(X^\top X)$), not to $1$.
  \yes{G} Standard use case: project to 2--3 dimensions for visualization while keeping maximal variance.
  \no{H} $U$ is orthogonal ($U^\top U=I$), not symmetric in general.
\end{answer}

\questiontitle{Picture-based interpretation question: PCA, ICA, and EM}
The original exam contained three picture-based subquestions.
State which interpretation is more likely in each case.
\begin{enumerate}
  \item PCA and ICA are applied to a 2D cloud of points, leading to outcomes A and B.
  Which picture is more likely the result of PCA and which one the result of ICA?
  \item Two plots show a dataset in a coordinate system spanned by principal components.
  Which is more likely: picture 1 is shown in $\mathrm{PC2}$ vs. $\mathrm{PC1}$ and picture 2 in $\mathrm{PC3}$ vs. $\mathrm{PC1}$, or vice versa?
  \item The left figure shows an original dataset, while figures $a$ and $b$ show results of running EM for a mixture of two Gaussians.
  Which is more likely: figure $a$ is after $5$ iterations and figure $b$ after $20$ iterations, or the other way around?
\end{enumerate}
\closequestion
\begin{answerprose}{PCA: A, ICA: B / picture 1 = PC2 vs.\ PC1 / a = 5 iterations, b = 20 iterations}
\begin{itemize}[leftmargin=1.7em,itemsep=1.5pt,topsep=2pt]
  \item \textbf{PCA vs.\ ICA:} PCA always returns \emph{orthogonal} axes aligned with the
        directions of maximal variance, so the picture with perpendicular axes is PCA (A).
        ICA finds statistically independent directions, which are in general \emph{not}
        orthogonal and follow the actual (possibly skewed) structure of the cloud (B).
  \item \textbf{PC order:} the projected variance decreases with the component index, so the
        scatter along PC2 must be larger than along PC3. The plot with the larger vertical
        spread is PC2 vs.\ PC1 (picture 1); the flatter one is PC3 vs.\ PC1 (picture 2).
  \item \textbf{EM iterations:} early in EM the two Gaussians still overlap and sit near their
        initialization; after more iterations they have converged onto the two clusters.
        Hence $a$ (poorly separated) is after $5$ iterations and $b$ (well fitted) after $20$.
\end{itemize}
\end{answerprose}

\questiontitle{PCA computation on a rank-1 matrix --- projected second sample}
\repeatnote

You are given four data samples, each of dimension $3$.
The data matrix reads
\[
X=\begin{pmatrix}
1&1&1\\
2&2&2\\
3&3&3\\
4&4&4
\end{pmatrix}.
\]
The eigenvectors of $X^\top X$ are
\[
u_1=\frac{1}{\sqrt3}(1,1,1)^\top,
\qquad
u_2=\frac{1}{\sqrt2}(-1,0,1)^\top,
\qquad
u_3=\frac{1}{\sqrt2}(-1,1,0)^\top.
\]
They are already normalized and already correctly ordered corresponding to eigenvalues in descending order.
Compute the PCA projection matrix $Y$ and enter the components of the projected second feature vector $y_2$.
Round to two decimal places.
\closequestion
\begin{answerprose}{(3.46, 0.00, 0.00)}
The projection is $Y=XU$, i.e.\ row-wise $y_i^{(k)}=x_i^\top u_k$. For the second sample
$x_2=(2,2,2)^\top$:
\[
y_2^{(1)}=\frac{2+2+2}{\sqrt3}=\frac{6}{\sqrt3}=2\sqrt3\approx\mathbf{3.46},
\]
\[
y_2^{(2)}=\frac{-2+0+2}{\sqrt2}=\mathbf{0},
\qquad
y_2^{(3)}=\frac{-2+2+0}{\sqrt2}=\mathbf{0}.
\]
Every row of $X$ is a multiple of $(1,1,1)^\top$ --- the matrix has rank $1$, so \emph{all}
variance lies along $u_1$ and the other components vanish for every sample.
\end{answerprose}

\questiontitle{PCA computation on a centered matrix --- projected third sample}
\repeatnote

You are given four data samples, each of dimension $3$.
The data matrix reads
\[
X=\begin{pmatrix}
-1&0&1\\
-2&0&2\\
-3&0&3\\
-4&0&4
\end{pmatrix}.
\]
The eigenvectors of the covariance matrix are
\[
u_1=\frac{1}{\sqrt2}(-1,0,1)^\top,
\qquad
u_2=\frac{1}{\sqrt2}(1,0,1)^\top,
\qquad
u_3=(0,1,0)^\top.
\]
They are already normalized and already correctly ordered corresponding to eigenvalues in descending order.
Compute the PCA projection matrix $Y$ and enter the components of the projected third sample $y_3$.
Round to two decimal places.
\closequestion
\begin{answerprose}{(4.24, 0.00, 0.00)}
With $Y=XU$ and $x_3=(-3,0,3)^\top$:
\[
y_3^{(1)}=\frac{(-3)(-1)+0+3\cdot1}{\sqrt2}=\frac{6}{\sqrt2}\approx\mathbf{4.24},
\]
\[
y_3^{(2)}=\frac{-3+0+3}{\sqrt2}=\mathbf{0},
\qquad
y_3^{(3)}=0\cdot(-3)+1\cdot0+0\cdot3=\mathbf{0}.
\]
All rows are multiples of $(-1,0,1)^\top$, so the data is one-dimensional along $u_1$.
\end{answerprose}

\questiontitle{PCA computation with a repeated sample --- projected third sample}
You are given four data samples, each of dimension $3$, collected in the data matrix
\[
X=\begin{pmatrix}
-4&0&-4\\
-2&0&2\\
3&0&3\\
3&0&3
\end{pmatrix}.
\]
Further assume that the eigenvectors of the covariance matrix of this data are
\[
u_1=\frac{1}{\sqrt2}(1,0,1)^\top,
\qquad
u_2=\frac{1}{\sqrt2}(-1,0,1)^\top,
\qquad
u_3=(0,1,0)^\top.
\]
They are already normalized and already correctly ordered corresponding to eigenvalues in descending order.
Compute the PCA projection matrix $Y$ and enter the components of the projected third sample $y_3$.
Round to two decimal places.
\closequestion
\begin{answerprose}{(4.24, 0.00, 0.00)}
With $Y=XU$ and $x_3=(3,0,3)^\top$:
\[
y_3^{(1)}=\frac{3\cdot1+0+3\cdot1}{\sqrt2}=\frac{6}{\sqrt2}\approx\mathbf{4.24},
\qquad
y_3^{(2)}=\frac{3\cdot(-1)+0+3\cdot1}{\sqrt2}=\mathbf{0},
\qquad
y_3^{(3)}=\mathbf{0}.
\]
The third sample lies exactly on the $(1,0,1)$ direction, so only the first component survives.
(The fourth sample is identical to the third, so it projects to the same point.)
\end{answerprose}

\questiontitle{Kernel PCA}
Which statements about Kernel PCA are correct?
Use the notation from the lecture: data points $x_1,\dots,x_n\in\mathcal X=\R^m$, kernel $k:\mathcal X\times\mathcal X\to\R$, feature map $\phi:\mathcal X\to\mathcal H$, feature covariance matrix
\[
C=\frac1n\sum_{i=1}^n\phi(x_i)\phi(x_i)^\top,
\]
and Gram matrix $K$ with entries $K_{ij}=\phi(x_i)^\top\phi(x_j)$.
Multiple answers are possible.
\begin{itemize}
  \choice A disadvantage of Kernel PCA is that $\phi$ always needs to be computed explicitly, which is computationally costly.
  \choice In general, any nonnegative and symmetric function $k:\mathcal X\times\mathcal X\to\R$ defines a kernel.
  \choice To compute the projection of a data point $x$ onto the eigenvectors of $K$, you need to center $x$ and $K$, as well as to normalize the eigenvectors of $K$.
  \choice The main task of Kernel PCA is to solve the eigenvalue problem of $C$ in feature space. This can be achieved by solving the eigenvalue problem of $K$ in a Euclidean space.
  \choice In general, $\phi$ is a nonlinear map that maps into a higher-dimensional feature space $\mathcal H$.
  \choice Linear PCA can be seen as a special case of Kernel PCA by using $k(x,y)=x^\top y$ as the associated kernel function.
  \choice For Kernel PCA, it is essential to center $K$ before computing its eigendecomposition, since the eigenvalues of $K$ can in general be negative, which leads to inconclusive results.
  \choice $K$ and $C$ always have the same dimensionality and the same rank.
\end{itemize}
\closequestion
\begin{answer}{C, D, E, F}
  \no{A} The kernel trick exists precisely so that $\phi$ is \emph{never} computed explicitly --- only inner products via $k$.
  \no{B} A kernel must be symmetric \emph{and positive semidefinite} (Mercer); nonnegativity of the function values is neither necessary nor sufficient.
  \yes{C} Centering (in feature space, applied to both $K$ and the new point) and normalizing the eigenvectors are both required for correct projections.
  \yes{D} The eigenproblem of $C$ in $\mathcal H$ reduces to the $n\times n$ eigenproblem of the Gram matrix $K$.
  \yes{E} Typical feature maps are nonlinear and map into a higher- (possibly infinite-) dimensional space.
  \yes{F} With the linear kernel $k(x,y)=x^\top y$, Kernel PCA reproduces ordinary PCA.
  \no{G} Tricky: centering \emph{is} needed, but the reason is wrong --- $K$ is PSD, so its eigenvalues are never negative. Centering removes the feature-space mean.
  \no{H} They share the same \emph{nonzero} eigenvalues (hence rank), but $K$ is $n\times n$ while $C$ lives in feature space --- the dimensionalities differ in general.
\end{answer}

\questiontitle{Independent Component Analysis --- version A}
Which statements about Independent Component Analysis (ICA) are correct?
Multiple answers are possible.
\begin{itemize}
  \choice Maximizing the negentropy maximizes the mutual information.
  \choice In general, ICA components are neither ranked nor orthogonal.
  \choice Non-Gaussianity can be measured by the second cumulant.
  \choice Criteria for independence are mutual information between components and non-Gaussianity of components.
  \choice The ICA solution is unique.
  \choice For ICA, one needs at least as many sources as observations.
  \choice Infomax maximizes the mutual information between components.
  \choice FastICA relies on whitening and rotating the data.
\end{itemize}
\closequestion
\begin{answer}{B, D, H}
  \no{A} Maximizing negentropy (non-Gaussianity) \emph{minimizes} the mutual information between components.
  \yes{B} Unlike PCA, the components carry no natural order and need not be orthogonal.
  \no{C} The second cumulant is the \emph{variance} --- it is blind to non-Gaussianity; one uses the 3rd/4th cumulants (skewness, kurtosis).
  \yes{D} Those are exactly the two standard independence criteria.
  \no{E} Only unique up to permutation, sign, and scaling of the components.
  \no{F} Reversed: one needs at least as many \emph{observations} as sources.
  \no{G} Infomax maximizes the output entropy and thereby \emph{minimizes} the mutual information between components.
  \yes{H} FastICA: whiten the data first, then rotate to maximize non-Gaussianity.
\end{answer}

\questiontitle{Independent Component Analysis --- version B}
Which statements about Independent Component Analysis (ICA) are correct?
Multiple answers are possible.
\begin{itemize}
  \choice Independence can be quantified by the mutual information between components or the non-Gaussianity of components.
  \choice Fast ICA uses whitening and rotating the data.
  \choice Maximizing the negentropy minimizes the mutual information.
  \choice For non-Gaussian sources, the ICA solution is unique up to permutation and scaling.
  \choice Non-Gaussianity can be measured by the second and third cumulant.
  \choice For ICA, the number of sources must be smaller than, or equal to, the number of observations.
  \choice ICA components are ranked and orthogonal.
  \choice Infomax maximizes the mutual information between components.
\end{itemize}
\closequestion
\begin{answer}{A, B, C, D, F}
  \yes{A} Mutual information and non-Gaussianity are the standard independence criteria.
  \yes{B} Whitening + rotation is the FastICA recipe.
  \yes{C} This is the correct direction: negentropy up $\Leftrightarrow$ mutual information down.
  \yes{D} Identifiability result for non-Gaussian sources: unique up to permutation and scaling (incl.\ sign).
  \no{E} Near-miss: the \emph{third} cumulant (skewness) does measure non-Gaussianity, but the \emph{second} is just the variance --- including it makes the statement false.
  \yes{F} $\#\text{sources}\le\#\text{observations}$ is the standard requirement.
  \no{G} Neither ranked nor orthogonal in general.
  \no{H} Infomax \emph{minimizes} the mutual information between components.
\end{answer}

\questiontitle{Independent Component Analysis --- version C}
\repeatnote

Which statements about Independent Component Analysis (ICA) are correct?
Multiple answers are possible.
\begin{itemize}
  \choice For ICA, one needs more sources than observations.
  \choice Maximizing the negentropy maximizes the mutual information.
  \choice Criteria for independence are mutual information between components and non-Gaussianity of components.
  \choice The ICA solution is unique.
  \choice Fast ICA relies on whitening and rotating the data.
  \choice Infomax maximizes the mutual information between components.
  \choice Non-Gaussianity can be measured by the second cumulant.
  \choice In general, ICA components are neither ranked nor orthogonal.
\end{itemize}
\closequestion
\begin{answer}{C, E, H}
  \no{A} Reversed: at least as many \emph{observations} as sources.
  \no{B} Maximizing negentropy \emph{minimizes} mutual information.
  \yes{C} The two standard independence criteria.
  \no{D} Unique only up to permutation, sign, and scaling.
  \yes{E} Whitening + rotation is the FastICA scheme.
  \no{F} Infomax minimizes the mutual information between components.
  \no{G} The second cumulant is the variance --- no information about Gaussianity.
  \yes{H} Correct --- no ranking, no orthogonality.
\end{answer}
\section{Factor Analysis}

\questiontitle{Factor Analysis: low-rank covariance part}
In factor analysis, part of the covariance function is formed by
\[
C_U = UU^\top = \sum_{i=1}^{\ell} u_i u_i^\top,
\]
where $U\in\R^{m\times \ell}$, the vectors $u_i$ are the columns of $U$, and $\ell<m$.
Which constraints does this impose on $C_U$?
Multiple answers are possible.
\begin{itemize}
  \choice The maximum rank of the matrix is $\ell$.
  \choice The number of zero eigenvalues of $C_U$ is greater or equal to $m-\ell$.
  \choice Any covariance matrix in $m$ dimensions can be written in the form above.
  \choice $C_U$ is singular.
\end{itemize}
\closequestion
\begin{answer}{A, B, D}
  \yes{A} $C_U$ is a sum of $\ell$ rank-1 terms, so $\mathrm{rank}(C_U)\le\ell$.
  \yes{B} A symmetric $m\times m$ matrix of rank $\le\ell$ has at least $m-\ell$ zero eigenvalues.
  \no{C} A full-rank covariance matrix (rank $m$) cannot be written with only $\ell<m$ columns --- that is exactly why FA adds the diagonal noise term $\Psi$.
  \yes{D} Rank $<m$ implies $\det C_U=0$, i.e.\ $C_U$ is singular.
\end{answer}

\questiontitle{Factor Analysis: dimensions of the objects}
We keep the notation from the lecture.
Here $n$ is the number of data samples, $m$ is the dimension of each sample, $\ell$ is the number of factors,
$y$ is the factor vector, $U$ is the factor loading matrix, $\Psi$ is the diagonal noise covariance matrix, and $\varepsilon$ is the noise vector.
Which statements are correct?
Multiple answers are possible.
\begin{itemize}
  \choice $\Psi$ has shape $m\times m$.
  \choice $\Psi$ has shape $\ell\times \ell$.
  \choice $\Psi$ has shape $n\times \ell$.
  \choice $U$ has shape $m\times \ell$.
  \choice $U$ has shape $m\times n$.
  \choice $U$ has shape $\ell\times m$.
  \choice $\varepsilon$ has dimension $m$.
  \choice $\varepsilon$ has dimension $\ell$.
  \choice $\varepsilon$ has dimension $n$.
  \choice $y$ has dimension $m$.
  \choice $y$ has dimension $\ell$.
  \choice $y$ has dimension $n$.
\end{itemize}
\closequestion
\begin{answerprose}{A, D, G, K}
Everything follows from the FA model equation
\[
x = Uy+\varepsilon,\qquad x\in\R^m,
\]
read off left to right: to produce an $m$-dimensional observation from an $\ell$-dimensional
factor vector, $U$ must be $m\times\ell$ (\textbf{D}); the factor vector $y$ has dimension
$\ell$ (\textbf{K}); the noise $\varepsilon$ is added to the observation, so it has dimension
$m$ (\textbf{G}); and its covariance $\Psi=\mathrm{Cov}(\varepsilon)$ is $m\times m$
(\textbf{A}). The number of samples $n$ never appears in any of these shapes --- every option
containing $n$ is automatically wrong.
\end{answerprose}

\questiontitle{Factor Analysis: compute the first row of the covariance matrix}
You are given the factor loading matrix
\[
U=\begin{pmatrix}
1&2\\
3&4\\
5&6
\end{pmatrix}
\]
and the diagonal noise covariance matrix
\[
\Psi=\begin{pmatrix}
5&0&0\\
0&5&0\\
0&0&5
\end{pmatrix}.
\]
Compute the first row of the covariance matrix
\[
C=UU^\top+\Psi.
\]
\closequestion
\begin{answerprose}{(10, 11, 17)}
The first row of $UU^\top$ is the first row of $U$, namely $(1,2)$, multiplied with each row of $U$:
\[
(UU^\top)_{11}=1\cdot1+2\cdot2=5,\qquad
(UU^\top)_{12}=1\cdot3+2\cdot4=11,\qquad
(UU^\top)_{13}=1\cdot5+2\cdot6=17.
\]
Adding $\Psi$ only changes the diagonal entry: $C_{11}=5+5=\mathbf{10}$, $C_{12}=\mathbf{11}$,
$C_{13}=\mathbf{17}$.
\end{answerprose}

\questiontitle{Factor Analysis: conceptual statements}
Which statements about Factor Analysis (FA) are correct?
Multiple answers are possible.
\begin{itemize}
  \choice The factors are unique only up to orthogonal transformations (rotations).
  \choice Correlations between observations can only be explained by factors since the noise covariance matrix $\Psi$ is diagonal.
  \choice FA describes the variability of observations in terms of latent variables and noise.
  \choice The factor loading matrix $U$ describes the independent part of the covariance of the observations.
  \choice FA decomposes the covariance matrix of the data as $C=UU^\top+\Psi$.
  \choice FA assumes Gaussian noise and non-Gaussian latent variables.
  \choice Factors explain the variance of the observations, while the correlation is explained by noise.
\end{itemize}
\closequestion
\begin{answer}{A, B, C, E}
  \yes{A} For any rotation $R$: $(UR)(UR)^\top=UU^\top$, so $U$ is only determined up to rotations.
  \yes{B} With diagonal $\Psi$, all off-diagonal covariance must come from $UU^\top$ --- i.e.\ from the factors.
  \yes{C} That is the FA model: shared latent factors plus independent noise.
  \no{D} Swapped: $UU^\top$ is the \emph{shared/correlated} part; the independent part is $\Psi$.
  \yes{E} Exactly the FA covariance decomposition.
  \no{F} Standard FA assumes \emph{both} Gaussian factors and Gaussian noise (that is what makes it a tractable linear-Gaussian model; ICA is the non-Gaussian relative).
  \no{G} Reversed: factors explain the \emph{correlations}, noise the independent leftover variance.
\end{answer}

\section{Clustering, Mixtures, and Projection Methods}

\questiontitle{$k$-means: one iteration, dataset A}
Assume you want to do $k$-means on the following 2-dimensional dataset:
\[
X=\{(1,0),(0,1),(-1,0),(0,-1),(2,0),(2,3),(5,0)\}.
\]
The initial cluster centers are $\mu=(0,1)$ and $\nu=(2,2)$.
Compute the coordinates of the new cluster centers $\mu^{\mathrm{new}}$ and $\nu^{\mathrm{new}}$ after one iteration of $k$-means.
\closequestion
\begin{answerprose}{$\mu^{\mathrm{new}}=(0,0)$, $\nu^{\mathrm{new}}=(3,1)$}
\textbf{Assignment step} (nearest center; squared distances suffice):
$(1,0),(0,1),(-1,0),(0,-1)$ are closer to $\mu=(0,1)$;
$(2,0),(2,3),(5,0)$ are closer to $\nu=(2,2)$
(e.g.\ $(2,0)$: $d^2_\mu=4+1=5$ vs.\ $d^2_\nu=0+4=4$).

\textbf{Update step} (means of the assigned points):
\[
\mu^{\mathrm{new}}=\tfrac14\bigl[(1,0)+(0,1)+(-1,0)+(0,-1)\bigr]=(\mathbf{0},\mathbf{0}),
\]
\[
\nu^{\mathrm{new}}=\tfrac13\bigl[(2,0)+(2,3)+(5,0)\bigr]=(\mathbf{3},\mathbf{1}).
\]
\end{answerprose}

\questiontitle{$k$-means and affinity propagation}
Which statements about $k$-means and affinity propagation are correct?
Multiple answers are possible.
\begin{itemize}
  \choice The $k$-means objective is to minimize the within-cluster sum of squared deviations.
  \choice In affinity propagation: the larger the preference, the smaller is the number of cluster centers.
  \choice In affinity propagation: the larger the damping factor, the faster the convergence.
  \choice In affinity propagation, cluster centers are data points.
  \choice In affinity propagation, the similarity between two points is typically chosen as their negative squared distance.
  \choice $k$-means always converges to the global optimum.
  \choice The parameter $k$ in $k$-means is automatically chosen by the algorithm to best fulfill the optimization objective.
  \choice $k$-means is sensitive to initialization of the cluster centers.
\end{itemize}
\closequestion
\begin{answer}{A, D, E, H}
  \yes{A} $k$-means minimizes $\sum_k\sum_{x\in C_k}\lVert x-\mu_k\rVert^2$.
  \no{B} Reversed: the preference is a point's a-priori suitability as an exemplar --- the larger it is, the \emph{more} cluster centers emerge.
  \no{C} Damping stabilizes the message updates against oscillations; a larger damping factor makes convergence \emph{slower}, not faster.
  \yes{D} Cluster centers (exemplars) are always actual data points.
  \yes{E} $s(i,k)=-\lVert x_i-x_k\rVert^2$ is the standard choice.
  \no{F} $k$-means only reaches a \emph{local} optimum of its objective.
  \no{G} $k$ is a user-chosen hyperparameter --- the algorithm never selects it.
  \yes{H} Different initial centers can lead to different local optima; that is why one restarts $k$-means several times.
\end{answer}

\questiontitle{$k$-means: one iteration, dataset B}
Assume you want to do $k$-means on the following 2-dimensional dataset:
\[
X=\{(1,0),(2,-1),(1,-2),(0,-1),(2,0),(2,1),(3,0),(3,1)\}.
\]
The initial cluster centers are $\mu=(0,-2)$ and $\nu=(4,1)$.
Compute the coordinates of the new cluster centers $\mu^{\mathrm{new}}$ and $\nu^{\mathrm{new}}$ after one iteration of $k$-means.
\closequestion
\begin{answerprose}{$\mu^{\mathrm{new}}=(1,-1)$, $\nu^{\mathrm{new}}=(2.5,0.5)$}
\textbf{Assignment:} $(1,0),(2,-1),(1,-2),(0,-1)\to\mu=(0,-2)$;
$(2,0),(2,1),(3,0),(3,1)\to\nu=(4,1)$
(e.g.\ $(2,0)$: $d^2_\mu=4+4=8$ vs.\ $d^2_\nu=4+1=5$).

\textbf{Update:}
\[
\mu^{\mathrm{new}}=\tfrac14\bigl[(1,0)+(2,-1)+(1,-2)+(0,-1)\bigr]=(\mathbf{1},\mathbf{-1}),
\]
\[
\nu^{\mathrm{new}}=\tfrac14\bigl[(2,0)+(2,1)+(3,0)+(3,1)\bigr]=(\mathbf{2.5},\mathbf{0.5}).
\]
\end{answerprose}

\questiontitle{$k$-means: one iteration, dataset C}
Assume you want to do $k$-means on the following 2-dimensional dataset:
\[
X=\{(1,-1),(0,0),(-1,-1),(0,-2),(2,0),(2,3),(5,0)\}.
\]
The initial cluster centers are $\mu=(0,0)$ and $\nu=(2,1)$.
Compute the coordinates of the new cluster centers $\mu^{\mathrm{new}}$ and $\nu^{\mathrm{new}}$ after one iteration of $k$-means.
\closequestion
\begin{answerprose}{$\mu^{\mathrm{new}}=(0,-1)$, $\nu^{\mathrm{new}}=(3,1)$}
\textbf{Assignment:} $(1,-1),(0,0),(-1,-1),(0,-2)\to\mu=(0,0)$;
$(2,0),(2,3),(5,0)\to\nu=(2,1)$
(e.g.\ $(2,0)$: $d^2_\mu=4$ vs.\ $d^2_\nu=1$).

\textbf{Update:}
\[
\mu^{\mathrm{new}}=\tfrac14\bigl[(1,-1)+(0,0)+(-1,-1)+(0,-2)\bigr]=(\mathbf{0},\mathbf{-1}),
\]
\[
\nu^{\mathrm{new}}=\tfrac13\bigl[(2,0)+(2,3)+(5,0)\bigr]=(\mathbf{3},\mathbf{1}).
\]
\end{answerprose}

\questiontitle{$k$-means: one iteration, dataset D}
Assume you want to do $k$-means on the following 2-dimensional dataset:
\[
X=\{(0,1),(1,2),(2,1),(1,0),(0,2),(-1,2),(0,3),(-1,3)\}.
\]
The initial cluster centers are $\mu=(2,0)$ and $\nu=(-1,4)$.
Compute the coordinates of the new cluster centers $\mu^{\mathrm{new}}$ and $\nu^{\mathrm{new}}$ after one iteration of $k$-means.
A sketch of the data points helps: the assignment step can largely be read off graphically.
Give the values with one decimal place.
\closequestion
\begin{answerprose}{$\mu^{\mathrm{new}}=(1.0,1.0)$, $\nu^{\mathrm{new}}=(-0.5,2.5)$}
\textbf{Assignment} (squared distances to $\mu=(2,0)$ vs.\ $\nu=(-1,4)$):
\[
\begin{array}{c|cc|c}
x & d^2_\mu & d^2_\nu & \to\\
\hline
(0,1) & 5 & 10 & \mu\\
(1,2) & 5 & 8 & \mu\\
(2,1) & 1 & 18 & \mu\\
(1,0) & 1 & 20 & \mu\\
(0,2) & 8 & 5 & \nu\\
(-1,2) & 13 & 4 & \nu\\
(0,3) & 13 & 2 & \nu\\
(-1,3) & 18 & 1 & \nu
\end{array}
\]
\textbf{Update:}
\[
\mu^{\mathrm{new}}=\tfrac14\bigl[(0,1)+(1,2)+(2,1)+(1,0)\bigr]=(\mathbf{1.0},\mathbf{1.0}),
\]
\[
\nu^{\mathrm{new}}=\tfrac14\bigl[(0,2)+(-1,2)+(0,3)+(-1,3)\bigr]=(\mathbf{-0.5},\mathbf{2.5}).
\]
\end{answerprose}

\questiontitle{$t$-SNE}
Which statements about $t$-distributed Stochastic Neighborhood Embedding are correct?
Multiple answers are possible.
\begin{itemize}
  \choice It models the neighborhood in both observation and embedding space by a Gaussian distribution.
  \choice It minimizes the Kullback--Leibler divergence between distributions in observation space and embedding space.
  \choice It models the neighborhood of a data point in observation space proportional to a Student's $t$-distribution.
  \choice It is usually trained by gradient descent.
\end{itemize}
\closequestion
\begin{answer}{B, D}
  \no{A} Gaussian in the \emph{observation} space only; the embedding space uses a Student's $t$-distribution (that is the ``$t$'' in $t$-SNE).
  \yes{B} The objective is $\KL{P}{Q}$ between the input-space and embedding-space similarity distributions.
  \no{C} Swapped: the $t$-distribution is used in the \emph{embedding} space; the observation space uses Gaussians.
  \yes{D} The KL objective is minimized by (stochastic) gradient descent.
\end{answer}

\questiontitle{Scaling and Projection Methods --- version A}
Which statements about Scaling and Projection Methods are correct?
Multiple answers are possible.
\begin{itemize}
  \choice Isomap uses a geodesic distance matrix, whose eigenvectors are the coordinates of the projected space.
  \choice The objective in $t$-SNE is minimization of the KL divergence between the similarities of the high-dimensional inputs and the low-dimensional mapped outputs.
  \choice In locally linear embedding (LLE), each data point is described as a linear combination of its neighbors.
  \choice The objective of scaling is to project data into a higher-dimensional space, e.g. to allow better model selection.
  \choice In self-organizing maps the objective is to minimize an energy (error) function.
  \choice Multidimensional scaling (MDS) keeps the distances between data points as well as possible.
  \choice LLE in general performs linear mappings to realize neighborhood-preserving embeddings.
  \choice In $t$-SNE, the similarity between data points $i$ and $j$ is given by their joint probability.
\end{itemize}
\closequestion
\begin{answer}{A, B, C, F, H}
  \yes{A} Isomap is classical MDS on geodesic distances; the leading eigenvectors of the (centered) distance matrix give the embedding coordinates.
  \yes{B} That is exactly the $t$-SNE objective.
  \yes{C} LLE reconstructs every point as a weighted linear combination of its neighbors and preserves those weights in the embedding.
  \no{D} Scaling/projection methods go to a \emph{lower}-dimensional space.
  \no{E} The Kohonen SOM update rules are \emph{not} the gradient of any global energy function --- there is no energy being minimized.
  \yes{F} Distance preservation is the definition of MDS.
  \no{G} LLE is a \emph{nonlinear} embedding method --- only the local reconstructions are linear.
  \yes{H} $t$-SNE uses symmetrized joint probabilities $p_{ij}$ and $q_{ij}$.
\end{answer}

\questiontitle{Projection and scaling methods --- version B}
Which statements about projection and scaling methods are correct?
Multiple answers are possible.
\begin{itemize}
  \choice Non-negative Matrix Factorization (NMF) performs a multiplicative decomposition of a positive semidefinite matrix into two positive semidefinite matrices.
  \choice Locally Linear Embedding (LLE) is a linear embedding method that aims to preserve the neighborhood of each data point.
  \choice Multidimensional Scaling (MDS) aims to find a low-dimensional representation of the data that preserves pairwise Euclidean distances.
  \choice Stochastic Neighbor Embedding (SNE) defines the similarity of two data points as a conditional probability while $t$-SNE defines it as a joint probability.
  \choice Isomap measures distances between data points as the length of the shortest path along a neighborhood graph.
\end{itemize}
\closequestion
\begin{answer}{C, D, E}
  \no{A} NMF factorizes an \emph{entrywise non-negative} matrix into entrywise non-negative factors --- positive semidefiniteness is a different concept and plays no role.
  \no{B} LLE preserves neighborhoods, but it is a \emph{nonlinear} method.
  \yes{C} That is the definition of (metric) MDS.
  \yes{D} SNE uses conditionals $p_{j\mid i}$; $t$-SNE symmetrizes them into joints $p_{ij}$.
  \yes{E} Geodesic distance $\approx$ shortest path through the $k$-nearest-neighbor graph.
\end{answer}

\questiontitle{Scaling and projection methods --- version C}
Tick the correct statements about scaling and projection methods.
We collect the original data points in a matrix $X$ and the mapped points in a matrix $Y$.
Multiple answers are possible.
\begin{itemize}
  \choice The $t$-distribution has heavier tails than the Gaussian distribution; thus in $t$-SNE the joint distributions $q_{ij}$ between $y_i$ and $y_j$ are computed using a $t$-distribution in order to alleviate the crowding problem.
  \choice Isomap tries to approximate the geodesic distance between two points on a manifold by constructing a neighborhood graph.
  \choice In SNE and $t$-SNE, the perplexity term is computed using the entropy of the similarities $q_{ij}$ of the mapped data points $y_i$ and $y_j$.
  \choice In SNE and $t$-SNE, the objective is the mean squared distance between $x_i$ and $y_i$.
  \choice Locally linear embedding and multidimensional scaling are probabilistic methods that rely on minimizing the KL divergence between the distributions of $X$ and $Y$.
  \choice In SNE and $t$-SNE, the computation of the similarity $p_{ij}$ between $x_i$ and $x_j$ relies on assuming a Gaussian distribution centered at $x_i$.
\end{itemize}
\closequestion
\begin{answer}{A, B, F}
  \yes{A} The heavy tails give moderately distant pairs more room in the embedding --- exactly the fix for the crowding problem.
  \yes{B} Isomap approximates geodesics by shortest paths in the neighborhood graph.
  \no{C} Tricky: perplexity is set on the \emph{input-space} similarities $p_{j\mid i}$ (to choose the Gaussian bandwidths), not on the $q_{ij}$ of the mapped points.
  \no{D} The objective is the KL divergence between similarity distributions, not a mean squared distance ($x_i$ and $y_i$ don't even live in the same space).
  \no{E} Neither LLE nor MDS is probabilistic or KL-based --- that describes SNE/$t$-SNE.
  \yes{F} $p_{j\mid i}$ comes from a Gaussian kernel centered at $x_i$.
\end{answer}

\questiontitle{Scaling and Projection Methods --- version D}
\repeatnote

Which statements about Scaling and Projection Methods are correct?
Multiple answers are possible.
\begin{itemize}
  \choice Multidimensional scaling (MDS) keeps the distances between data points as well as possible.
  \choice The objective of scaling is to project data into a higher-dimensional space, e.g. to allow better model selection.
  \choice Isomap uses a geodesic distance matrix, whose eigenvectors are the coordinates of the projected space.
  \choice In $t$-SNE, the similarity between data points $i$ and $j$ is given by their joint probability.
  \choice In self-organizing maps the objective is to minimize an energy (error) function.
  \choice The objective in $t$-SNE is minimization of the KL divergence between the similarities of the high-dimensional inputs and the low-dimensional mapped outputs.
  \choice In locally linear embedding (LLE), each data point is described as a linear combination of its neighbors.
  \choice LLE in general performs linear mappings to realize neighborhood-preserving embeddings.
\end{itemize}
\closequestion
\begin{answer}{A, C, D, F, G}
  \yes{A} Distance preservation is the definition of MDS.
  \no{B} \emph{Lower}-dimensional, not higher.
  \yes{C} Isomap = MDS on the geodesic distance matrix; its eigenvectors give the coordinates.
  \yes{D} $t$-SNE uses joint probabilities $p_{ij}$, $q_{ij}$.
  \no{E} SOMs have no global energy function.
  \yes{F} That is the $t$-SNE objective.
  \yes{G} LLE's local reconstruction property.
  \no{H} LLE is nonlinear overall --- only the local reconstructions are linear.
\end{answer}

\questiontitle{Scaling and Projection Methods --- version E}
Decide which statements about Scaling and Projection Methods are correct.
Multiple answers are possible.
\begin{itemize}
  \choice The objective of scaling and projection methods is to project data into a lower-dimensional space, e.g. to allow better model selection.
  \choice In self-organizing maps the objective is to minimize an energy (error) function.
  \choice Projection pursuit performs projections such that the extracted signal is as Gaussian as possible.
  \choice Generative Topographic Mapping (GTM) is a non-linear latent variable model which maps latent variables to observations.
  \choice Isomap uses a geodesic distance matrix, whose largest eigenvectors are the coordinates of the projected space.
  \choice The objective in $t$-SNE is minimization of the KL divergence between the similarities of the (high-dimensional) inputs and the (low-dimensional) mapped outputs.
  \choice In locally linear embedding (LLE), each data point is described as a linear combination of its neighbors.
  \choice Multidimensional Scaling (MDS) projects data to a higher-dimensional space in order to increase the distance between data points.
\end{itemize}
\closequestion
\begin{answer}{A, D, E, F, G}
  \yes{A} This is the correct version of the recurring distractor: the projection goes into a \emph{lower}-dimensional space.
  \no{B} The Kohonen SOM update rules are not the gradient of any energy function --- nothing is being minimized.
  \no{C} Reversed: projection pursuit looks for maximally \emph{non}-Gaussian (``interesting'') projections --- a Gaussian projection is the least informative one.
  \yes{D} GTM maps a grid of latent points through a nonlinear function into data space --- a nonlinear latent variable model.
  \yes{E} The eigenvectors belonging to the largest eigenvalues of the (centered) geodesic distance matrix give the Isomap coordinates.
  \yes{F} The $t$-SNE objective.
  \yes{G} LLE's defining property.
  \no{H} MDS embeds into a \emph{lower}-dimensional space while \emph{preserving} distances --- it neither adds dimensions nor inflates distances.
\end{answer}
\questiontitle{Non-negative Matrix Factorization (NMF)}
As in the lecture, let
\[
X=YU^\top,
\qquad Y\in\R^{n\times \ell},
\qquad U\in\R^{m\times \ell},
\]
with non-negativity constraints.
Which statements about NMF are correct?
Multiple answers are possible.
\begin{itemize}
  \choice NMF objectives are the KL divergence and the Euclidean distance.
  \choice NMF is always exactly solvable.
  \choice $X\in\R^{n\times m}$.
  \choice The non-negativity constraints make the representation of the observations purely additive.
  \choice The outer product representation reads $X=\sum_{k=1}^{\ell} y_k u_k^\top$.
  \choice NMF factorizes any positive or negative definite matrix into two positive definite matrices.
\end{itemize}
\closequestion
\begin{answer}{A, C, D, E}
  \yes{A} The two standard NMF objectives are the squared Euclidean distance and the (generalized) KL divergence.
  \no{B} Exact NMF is NP-hard in general; the multiplicative update rules only find local optima of the chosen objective.
  \yes{C} $Y\in\R^{n\times\ell}$ times $U^\top\in\R^{\ell\times m}$ gives $X\in\R^{n\times m}$.
  \yes{D} Since no negative entries exist, components can never cancel each other --- the representation is purely additive (parts-based).
  \yes{E} Column-times-row expansion of the product: $X=\sum_k y_k u_k^\top$.
  \no{F} Nonsense --- definiteness plays no role in NMF, and a negative definite matrix cannot even have a non-negative factorization.
\end{answer}

\questiontitle{Mixturetown bread question}
In Mixturetown there are two bakers, Albert $(A)$ and Bibi $(B)$.
Their loaf weights are modeled as Gaussian distributions
\[
f(x;\mu_A,\sigma_A),\qquad f(x;\mu_B,\sigma_B),
\]
with
\[
\mu_A=1\text{ kg},\quad \sigma_A=0.1\text{ kg},\qquad
\mu_B=1.5\text{ kg},\quad \sigma_B=0.2\text{ kg}.
\]
Seventy-five percent of all loaves are made by Bibi.
You observe a loaf of weight $1.1$ kg.
Determine the probability that Bibi made this loaf.
Use the given density values
\[
f(1.1;1,0.1)=2.41971,
\qquad
f(1.1;1.5,0.2)=0.26995.
\]
Round to two decimal places.
\closequestion
\begin{answerprose}{0.25}
Bayes' rule with priors $P(A)=0.25$, $P(B)=0.75$:
\[
P(B\mid x=1.1)
=\frac{P(B)\,f(1.1;\mu_B,\sigma_B)}{P(A)\,f(1.1;\mu_A,\sigma_A)+P(B)\,f(1.1;\mu_B,\sigma_B)}
\]
\[
=\frac{0.75\cdot 0.26995}{0.25\cdot 2.41971+0.75\cdot 0.26995}
=\frac{0.20246}{0.60493+0.20246}
=\frac{0.20246}{0.80739}
\approx\mathbf{0.25}.
\]
Despite Bibi's large prior, the loaf weight is much more typical of Albert's ovens.
\end{answerprose}

\questiontitle{Mixture model posterior probability}
\repeatnote

Assume the following case:
\begin{itemize}
  \item The prior probability of component $j$ is $0.1$.
  \item The probability for observing a data sample $x$ given the total parameter set $\theta$ is $0.2$.
  \item The likelihood for observing a data sample $x$, given component $j$ and corresponding parameter vector, is $0.5$.
\end{itemize}
Compute the posterior probability of component $j$.
Round to two decimal places.
\closequestion
\begin{answerprose}{0.25}
Bayes' rule for mixtures:
\[
P(j\mid x)=\frac{P(j)\,p(x\mid j)}{p(x\mid\theta)}=\frac{0.1\cdot 0.5}{0.2}=\mathbf{0.25}.
\]
\end{answerprose}

\questiontitle{Mixture model: infer the likelihood from posterior, prior, and evidence}
\repeatnote

Assume the following case:
\begin{itemize}
  \item The prior probability of component $j$ is $0.1$.
  \item The posterior probability of component $j$ is $0.3$.
  \item The probability for observing a data sample $x$ given the total parameter set $\theta$ is $0.2$.
\end{itemize}
Compute the likelihood for observing a data sample $x$ given component $j$ and the corresponding parameter vector.
Round to two decimal places.
\closequestion
\begin{answerprose}{0.60}
Rearrange Bayes' rule $P(j\mid x)=\dfrac{P(j)\,p(x\mid j)}{p(x\mid\theta)}$ for the likelihood:
\[
p(x\mid j)=\frac{P(j\mid x)\;p(x\mid\theta)}{P(j)}=\frac{0.3\cdot 0.2}{0.1}=\mathbf{0.60}.
\]
\end{answerprose}

\questiontitle{Mixture of Gaussians: posterior of the first Gaussian}
The data generation process in a mixture of Gaussians model is as follows:
first choose one Gaussian according to its weight in the mixture, then draw a sample from that Gaussian.
You have fitted a mixture of two Gaussians to 1D data:
\[
p(x)=\alpha_1 f(x;\mu_1,\sigma_1)+\alpha_2 f(x;\mu_2,\sigma_2),
\]
with
\[
\alpha_1=0.3,\quad \alpha_2=0.7,\quad \mu_1=0,\quad \sigma_1=1,\quad \mu_2=3,\quad \sigma_2=1.
\]
Relevant densities are
\[
f(1;0,1)=0.242,
\qquad
f(1;3,1)=0.054.
\]
Calculate the probability that the datapoint $x=1$ has been produced by the first Gaussian.
Round to two decimal places.
\closequestion
\begin{answerprose}{0.66}
\[
P(G_1\mid x=1)
=\frac{\alpha_1 f(1;0,1)}{\alpha_1 f(1;0,1)+\alpha_2 f(1;3,1)}
=\frac{0.3\cdot 0.242}{0.3\cdot 0.242+0.7\cdot 0.054}
\]
\[
=\frac{0.0726}{0.0726+0.0378}
=\frac{0.0726}{0.1104}
\approx\mathbf{0.66}.
\]
\end{answerprose}

\questiontitle{KL divergence for continuous uniform distributions --- variant 2}
Assume $Y$ and $Z$ are uniformly distributed on $[0,\theta_Y]$ and $[0,\theta_Z]$, respectively.
For $\theta_Y=e^4$ and $\theta_Z=e^{10}$, compute
\[
D_{\mathrm{KL}}(p_Y\|p_Z)=\int_{-\infty}^{\infty} p_Y(x)\log\!\left(\frac{p_Y(x)}{p_Z(x)}\right)\dd x.
\]
Use the natural logarithm and round to two decimal places.
\closequestion
\begin{answerprose}{6.00}
Same derivation as before:
\[
D_{\mathrm{KL}}(p_Y\|p_Z)=\ln\frac{\theta_Z}{\theta_Y}=\ln\frac{e^{10}}{e^{4}}=10-4=\mathbf{6.00}.
\]
\end{answerprose}

\questiontitle{KL divergence for two discrete distributions}
\repeatnote

You are given two probability distributions $p$ and $q$:
\[
p(x=0)=0.8,\quad p(x=1)=0.2,
\qquad
q(x=0)=0.5,\quad q(x=1)=0.5.
\]
Compute the KL divergence $D_{\mathrm{KL}}(p\|q)$.
Use the natural logarithm and round to two decimal places.
\closequestion
\begin{answerprose}{0.19}
\[
D_{\mathrm{KL}}(p\|q)=\sum_x p(x)\ln\frac{p(x)}{q(x)}
=0.8\ln\frac{0.8}{0.5}+0.2\ln\frac{0.2}{0.5}
=0.8\cdot 0.470+0.2\cdot(-0.916)
\]
\[
=0.376-0.183\approx\mathbf{0.19}.
\]
\end{answerprose}

\questiontitle{Clustering --- version A}
Which statements about Clustering are correct?
Multiple answers are possible.
\begin{itemize}
  \choice In affinity propagation, the cluster centers are always data points.
  \choice Clustering (or Cluster Analysis) is both a regression and a classification technique.
  \choice In both $k$-means and fuzzy $k$-means, the membership distribution is always given by a binary random variable.
  \choice In mixture models, an observation is assigned to the component with the largest prior distribution.
  \choice Top-down clustering relies on building the minimal spanning tree and then in each step removes the largest of the remaining edges.
  \choice In a mixture of Gaussians, the maximum likelihood and the maximum posterior always have the same value.
  \choice The agglomerative hierarchical clustering algorithm repeatedly fuses the two maximally distant clusters.
  \choice Similarity-based clustering does not require objects to be represented via feature vectors.
\end{itemize}
\closequestion
\begin{answer}{A, H}
  \yes{A} Affinity propagation selects \emph{exemplars} --- actual data points --- as cluster centers.
  \no{B} Clustering is unsupervised --- it is neither regression nor classification (both are supervised).
  \no{C} \emph{Fuzzy} $k$-means is exactly the variant with continuous (soft) memberships; only plain $k$-means is binary.
  \no{D} Assignment goes to the largest \emph{posterior} (responsibility), which also accounts for how well the component explains the observation --- not the prior alone.
  \no{E} Removing the longest edges of a minimal spanning tree reproduces \emph{single-linkage} (bottom-up) clusters; top-down (divisive) clustering instead splits clusters recursively.
  \no{F} MAP includes the prior term, so the two optima and their values differ in general.
  \no{G} Agglomerative clustering fuses the two \emph{closest} clusters.
  \yes{H} It only needs pairwise similarities --- objects never have to live in a vector space.
\end{answer}

\questiontitle{Clustering --- version B}
Which statements about Clustering are correct?
Multiple answers are possible.
\begin{itemize}
  \choice $k$-means is faster than MoG and not sensitive to initialization.
  \choice Both $k$-means and MoG assign an observation to a single cluster only.
  \choice The mixing coefficients in a Mixture of Gaussians model are represented by the posteriors.
  \choice The maximum a posteriori estimation is a way to avoid near-zero variances and singular solutions in the maximum likelihood estimation of mixture models.
  \choice Given an observation $x$, a component $j$, and a parameter vector $\theta$, the posterior $p(j\mid x,\theta)$ describes how likely it is that $x$ was produced by $j$.
  \choice The EM algorithm for MoG tries to maximize the likelihood with respect to the cluster centers, the covariance matrices, and the mixing coefficients.
  \choice The $k$-means algorithm is a simplification of the Mixture of Gaussians model.
  \choice Clustering is an unsupervised technique that aims at finding separate regions of high data density called clusters.
\end{itemize}
\closequestion
\begin{answer}{D, E, F, G, H}
  \no{A} Half-true trap: $k$-means \emph{is} faster than MoG, but it is definitely sensitive to initialization.
  \no{B} $k$-means does, but MoG assigns \emph{soft} (probabilistic) memberships to all components.
  \no{C} The mixing coefficients are the \emph{priors} $\alpha_j$, not the posteriors.
  \yes{D} The prior in MAP regularizes the covariances away from collapsing onto single points (singular ML solutions).
  \yes{E} That is the definition of the responsibility $p(j\mid x,\theta)$.
  \yes{F} Means, covariances, and mixing coefficients are exactly the parameters EM optimizes.
  \yes{G} $k$-means is the hard-assignment, equal-spherical-covariance limit of MoG.
  \yes{H} Standard definition of clustering.
\end{answer}

\questiontitle{Hierarchical clustering: single linkage vs. complete linkage}
You are given $8$ datapoints in $\R^2$:
\[
\{(0,1),(1,1),(2,1),(8,1),(0,-1),(1,-1),(2,-2),(8,-1)\}.
\]
Agglomerative hierarchical clustering is run until only two clusters remain, once with single linkage and once with complete linkage.
Two resulting clusterings are shown in the exam figure.
Check the correct statements.
\begin{itemize}
  \choice Clustering 1 has been obtained by single linkage.
  \choice Clustering 2 has been obtained by single linkage.
  \choice Clustering 2 has been obtained by complete linkage.
  \choice Clustering 1 has been obtained by complete linkage.
  \choice One of the two clusterings could theoretically have been the result of either linkage function.
\end{itemize}
\closequestion
\begin{answer}{B, D, E}
  \no{A} See D.
  \yes{B} Single linkage merges via the \emph{closest} pair, so it chains the six left points into one elongated cluster and leaves the two distant points at $x=8$ as the second cluster --- the chained partition is clustering 2.
  \no{C} See B.
  \yes{D} Complete linkage merges via the \emph{farthest} pair, so it penalizes large cluster diameters and produces the compact partition --- clustering 1.
  \yes{E} Correct: the linkage functions only differ once cluster diameters start to matter; for this dataset one of the two shown partitions can arise under either criterion.
\end{answer}

\questiontitle{Hierarchical clustering with average linkage on letters}
\repeatnote

You are given the set of data points
\[
x=\{a,c,e,h,i,m,n,s\},
\]
where the similarity between two points $s(x,x')$ is defined as the distance between the ordered letters in the alphabet,
for example $s(a,b)=1$, $s(a,c)=2$, and so on.
Perform agglomerative clustering using average linkage,
\[
d_{\mathrm{avg}}(A,B)=\frac{1}{n_A}\frac{1}{n_B}\sum_{a\in A}\sum_{b\in B}s(a,b),
\]
and stop when the distance between clusters is greater than $7$.
How many clusters are formed in total, what is the maximum number of data points in any one cluster, and what is the minimum distance between any two of the formed clusters?
\closequestion
\begin{answerprose}{3 clusters / max size 4 / min distance 8}
Alphabet positions: $a=1$, $c=3$, $e=5$, $h=8$, $i=9$, $m=13$, $n=14$, $s=19$.
Merging the closest clusters step by step:
\begin{enumerate}[leftmargin=2.2em,itemsep=1.5pt,topsep=2pt]
  \item $\{h,i\}$ and $\{m,n\}$ at distance $1$ each.
  \item $\{a,c\}$ at distance $2$.
  \item $\{a,c\}+\{e\}$: $d_{\mathrm{avg}}=\tfrac{4+2}{2}=3$.
  \item $\{h,i\}+\{m,n\}$: $d_{\mathrm{avg}}=\tfrac{5+6+4+5}{4}=5$.
  \item Remaining distances:
        $d(\{a,c,e\},\{h,i,m,n\})=\tfrac{7+8+12+13+5+6+10+11+3+4+8+9}{12}=8$ and
        $d(\{h,i,m,n\},\{s\})=\tfrac{11+10+6+5}{4}=8$ --- both $>7$, so we \textbf{stop}.
\end{enumerate}
Result: \textbf{3 clusters} $\{a,c,e\}$, $\{h,i,m,n\}$, $\{s\}$; the largest contains
\textbf{4} points; the minimum distance between the formed clusters is \textbf{8}.
\end{answerprose}

\questiontitle{Biclustering}
Which statements about Biclustering are correct?
Multiple answers are possible.
\begin{itemize}
  \choice A common biclustering method is the variance maximization method.
  \choice FABIA is a generative model.
  \choice FABIA defines biclusters as inner products.
  \choice In FABIA, the likelihood is analytically intractable and requires, for example, a variational EM algorithm.
  \choice Row and column elements must belong to exactly one bicluster or to no bicluster at all.
  \choice Biclustering simultaneously clusters the rows and columns of a matrix.
\end{itemize}
\closequestion
\begin{answer}{B, D, F}
  \no{A} Made-up name --- not a biclustering method.
  \yes{B} FABIA is a generative model of the data matrix (factor analysis with sparse priors).
  \no{C} Near-miss: FABIA models a bicluster as an \emph{outer} product $\lambda z^\top$ of a sparse loading and a sparse factor, not an inner product.
  \yes{D} The FABIA likelihood has no closed form; it is trained with variational EM.
  \no{E} Biclusters may overlap --- rows and columns can belong to several biclusters or to none.
  \yes{F} That is the definition of biclustering.
\end{answer}

\end{document}
